\begin{abstract}
We compare eight language models with prediction-market prices on \polymarket and
\kalshi. The benchmark has 10{,}536 market moments, and the headline results use
the 4{,}865 held-out primary rows. The models never see a price or order book,
and most rows carry no news, so they work with less information than the market.
The market beats every model. The gap is in sorting, not calibration: the models
separate likely from unlikely outcomes less well than the crowd does. Resolution
accounts for 79 to 99 per cent of each model's shortfall, and no monotone
rescaling recovers more than about a fifth of it, or two fifths on the non-sports
rows, where the gap is widest. News improves sorting for nearly every model,
significantly only on the window whose news carries a leakage caveat, and never
to the market's level. The models are also wrong together. On rows the market got
right, two models land on the same wrong side 2.7 to 4.3 times as often as
independent forecasters would. Simulated forecasters that share only each
model's fitted line to the price reach 1.4 to 2.2. Adding each row's features or
its event still leaves them below the real rate. News reduces this agreement only
slightly, and we cannot tell missing information apart from shared training or
the shared prompt. No combination of the models, equal-weighted or learned,
matches the market. Returns are no worse in thin markets, though the accuracy
gap narrows in deeper order books, clearly on one window and less robustly on
another. These results are exploratory.
We release the benchmark, model outputs and code, except one window's market
data, which cannot be redistributed.
\end{abstract}
